\section{\'Etale coverings of a normal scheme}
\label{I.10}

\begin{proposition}
\label{I.10.1}
Let $X$ be an \'etale separated prescheme over a normal
connected $Y$ with field~$K$. Then the
connected components $X_i$ of~$X$ are integral, their
fields $K_i$ are finite separable extensions of~$K$,
and $X_i$ is identified with a nonempty open subset of
the normalization of~$X$ in $K_i$ (hence $X$ with a
dense open subset of the normalization of~$Y$ in
$R(X)=L=\prod K_i$).
\end{proposition}

By~\ref{I.9.10}, $X$ is normal; a fortiori its local
rings are integral domains, so the connected components
of~$X$ are irreducible. Since $X_i$ is normal, and finite
and dominant over~$Y$, it follows from a special case
(moreover, almost trivial) of the Main Theorem that
$X_i$ is an open subset of the normalization of~$X$
in the field $K_i$ of~$X$.

\begin{corollary}
\label{I.10.2}
Under the conditions of~\ref{I.10.1}, $X$ is finite
over~$Y$ (\ie an \'etale covering of~$Y$) if and only
if $X$ is isomorphic to the normalization $Y'$ of~$Y$
in $L=R(X)$ (the ring of rational functions on~$X$).
\end{corollary}

Indeed, one knows that this normalization is finite
over~$Y$ ($Y$ being normal and $R/K$ separable);
conversely, if $X$ is finite over~$Y$, it is finite
over~$Y'$, so its image in $Y'$ is closed, and on
the other hand it is dense.

An algebra $L$ of finite rank over~$K$ will be called
\emph{unramified over~$X$} (or simply unramified
over~$K$, if $X$ is understood) if $L$ is a separable
algebra over~$K$, \ie a direct product of separable
extensions $K_i$, and if the normalization $Y'$ of~$Y$
in $L$ (the disjoint sum of the normalizations of~$Y$
in the $K_i$) is unramified ($=$ \'etale by~\ref{I.9.11})
over~$Y$. Thus:

\begin{corollary}
\label{I.10.3}
For every $X$ finite over~$Y$ whose every irreducible
component dominates $Y$, let $R(X)$ be the ring of
rational functions on~$X$ (the product of the local
% original p. 22
rings of the generic points of the irreducible
components of~$X$), so that $X\mto R(X)$ is a functor
with values in the algebras of finite rank over
$K=R(Y)$. This functor establishes an equivalence of
the category of connected \'etale coverings of~$Y$
with the category of extensions $L$ of $K$ unramified
over~$Y$.
\end{corollary}

The inverse functor is the normalization functor.

Suppose that $Y$ is affine, hence defined by a normal
ring $A$ with field of fractions~$K$. Let $L$ be a
finite extension of~$K$ which is a direct product of
fields; then, by definition, the normalization $Y'$
of $Y$ in $L$ is isomorphic to $\Spec(A')$, where $A'$
is the normalization of~$A$ in~$L$. To say that $L$
is unramified over~$Y$ means that $A'$ is unramified
(or equivalently, \'etale) over~$A$. If $A$ is local,
this amounts to saying that the local rings
$A^\prime_{\goth{n}}$ (where $\goth{n}$ runs through
the finite set of maximal ideals of~$A'$, \ie its
prime ideals inducing the maximal ideal $\goth{m}$
of~$A$) are unramified ($=$ \'etale) over the local
ring~$A$. Finally, note also that the discriminant
criterion (\ref{I.4.10}) can also be applied in this
situation (more generally, a variant of the said
criterion should be stated as follows, without a
preliminary flatness condition when $X$ dominates
$Y$, with $Y$ nevertheless assumed locally integral:
$A\to B$ and $B\to B\otimes_A K=L$ are injective---then
$\trace_{L/K}$ is defined---and $\trace_{L/K}(xy)$
induces a \emph{fundamental bilinear form}
$B\times B\to A$, \ie there exist $x_i\in B$
($1\le i\le n$, $n=$ rank of~$L$ over~$K$) such that
$\trace(x_ix_j)\in A$ for every $i$,~$j$ and
$\det(\trace(x_ix_j))$ is invertible in~$A$).

The sorites (4.6) immediately implies the sorites
of unramifiedness in the classical setting:

\begin{proposition}
\label{I.10.4}
Let $Y$ be a normal integral prescheme, with field~$K$.
\textup{(i)} $K$ is unramified over~$Y$.
\textup{(ii)} If $L$ is an extension of~$K$ unramified
over~$Y$, if $Y'$ is a normal prescheme with field $L$
dominating $Y$ (for example, the normalization of~$Y$
in~$L$), and $M$ an extension of~$L$ unramified
over~$Y'$, then $M/K$ is unramified over~$X$
(\emph{transitivity} of unramifiedness).
\textup{(iii)} Let $Y'$ be a normal integral prescheme
dominating $Y$, with field $K'/K$; if $L$ is an
extension of~$K$ unramified over~$Y$, then
$L\otimes_K K'$ is an extension of~$K'$ unramified
over~$Y'$ (the \emph{translation} property).
\end{proposition}

Moreover:
% original p. 23

\begin{corollary}
\label{I.10.5}
Under the conditions of \textup{(iii)}, if
$Y=\Spec(A)$, $Y'=\Spec(A')$, then the normalization
$\bar{A'}$ of~$Y'$ in $L'=L\otimes_K K'$ is identified
with $\bar{A}\otimes_A A'$, where $\bar{A}$ is the
normalization of $A$ in $L$.
\end{corollary}

Usually, people (who are reluctant to consider rings
which are not integral domains, even if they are
direct products of fields) state the translation
property in the following (weaker) form:

\begin{corollary}
\label{I.10.6}
Under the conditions of \textup{(iii)}, let $L_1$ be
a \emph{compositum} of~$L/K$ (unramified over~$Y$)
and~$K'/K$. Then $L_1/K'$ is unramified over~$Y'$.
When $Y=\Spec(A)$, $Y'=\Spec(A')$, one moreover has
\[
\overline{A'}=A[\overline{A},A']
\]
\ie the ring $\overline{A'}$, the normalization
of~$A'$ in $L_1$, is the $A$-algebra generated by
$A'$ and the normalization $\overline{A}$ of~$A$
in $L$.
\end{corollary}

This last fact is false without the unramifiedness
assumption, even in the case of composita of number
fields\ldots

To conclude this no., we shall give the interpretation
of the notion of \'etale covering corresponding to
its intuitive picture: there should be the ``maximum
number'' of points above the point $y\in Y$ under
consideration, and in particular there should not be
``several coincident points'' above~$y$. To prove
results in this direction with all the desired
generality, we shall admit here Proposition~\ref{I.10.7}
below (whose proof will appear in the multiplodoque,
Chap\ptbl IV, par\ptbl 15, and uses Chevalley's
technique of constructible sets, and a little descent
theory\ldots).

A morphism of finite type $f\colon X\to Y$ is said
to be \emph{universally open} if for every extension
of the base $Y'\to Y$ (with $Y'$ locally noetherian)
the morphism $f'\colon X'=X\times_Y Y'\to Y'$ is
open, \ie sends open sets to open sets. One can,
moreover, restrict oneself to the case where $Y'$
is of finite type over~$Y$ (and even where $Y'$ is
of the form $Y[t_1,\dots,t_r]$, with the $t_i$
indeterminates). A universally open morphism is
a fortiori open (the converse being false);
on the other hand, if $f$ is open, with $X$ and $Y$
irreducible, then all the components of all the
fibers of~$f$ have the same dimension (namely, the
dimension of the generic fiber
% original p. 24
$f^{-1}(z)$, with $z$ the generic point of~$Y$).
Finally, if $Y$ is normal, this last condition
already implies that $f$ is \emph{universally} open
(Chevalley's theorem). It follows, for example,
that if $f\colon X\to Y$ is a \emph{quasi-finite}
morphism, with $Y$ normal and irreducible, then $f$
is universally open (or equivalently, open) if and
only if every irreducible component of~$X$ dominates~$Y$.
Recall also that a flat morphism (of finite type),
being open, is also universally open. With these
preliminaries in place, let us ``recall'' the following

\begin{proposition}
\label{I.10.7}
Let $f\colon X\to Y$ be a quasi-finite separated
universally open morphism. For every $y\in Y$, let
$n(y)$ be the ``geometric number of points of the
fiber $f^{-1}(y)$'', equal to the sum of the separable
degrees of the residue field extensions
$\kres(x)/\kres(y)$, for the points $x\in f^{-1}(y)$.
Then the function $y\mto n(y)$ on~$Y$ is upper
semicontinuous. For it to be constant in a neighborhood
of the point $y$ (\ie for one to have $n(y)=n(z_i)$,
where the $z_i$ are the generic points of the
irreducible components of~$Y$ containing $y$), it is
necessary that there exist a neighborhood $U$ of~$y$
such that $X|U$ is \emph{finite} over~$U$.\footnote{\Cf
EGA IV 15.5.1.}
\end{proposition}

\begin{corollary}
\label{I.10.8}
If $y\mto n(y)$ is constant and $Y$ is geometrically
unibranch\footnote{For the definition, \cf below,
\No\ref{I.11}, p\ptbl\pageref{I.11}.}, the irreducible
components of~$X$ are disjoint.
\end{corollary}

\begin{proposition}
\label{I.10.9}
Let $f\colon X\to Y$ be an \emph{\'etale} separated
morphism. With the notation of~\ref{I.10.7}, the
function $y\mto n(y)$ is upper semicontinuous. For
it to be constant in a neighborhood of the point $y$
(\ie for one to have $n(y)=n(z_i)$, where the $z_i$
are the generic points of the irreducible components
of~$Y$ containing $y$), it is necessary and sufficient
that there exist an open neighborhood $U$ of~$y$
such that $X|U$ is finite over U, \ie is an
\emph{\'etale covering} of~$U$.
\end{proposition}

\begin{corollary}
\label{I.10.10}
For a separated \'etale morphism $f\colon X\to Y$,
with $Y$ connected, to be finite (\ie to make $X$
an \emph{\'etale covering} of~$Y$), it is necessary
and sufficient that all the fibers of~$f$ have the
same geometric number of points.
\end{corollary}

In~\ref{I.10.7} and its corollary, there was no
normality assumption on~$Y$. If one makes such an
assumption, one obtains the stronger statement
(most often taken as the definition of a covering
being unramified):

\begin{theorem}
\label{I.10.11}
Let
% original p. 25
$f\colon X\to Y$ be a \emph{quasi-finite} separated
morphism. Suppose that~$Y$ is \emph{irreducible},
that every component of~$X$ dominates $Y$, and that
$X$ is reduced (\ie $\cal{O}_X$ has no nilpotent
elements). Let $n$ be the degree of~$X$ over~$Y$
(the sum of the degrees, over the field $K$ of~$Y$,
of the fields $K_i$ of the irreducible components
$X_i$ of~$X$). Let $y$ be a normal point of~$Y$.
Then the geometric number $n(y)$ of points of~$X$
above~$y$ is $\le n$, with equality if and only if
there exists an open neighborhood $U$ of~$y$ such
that $X|U$ is an \emph{\'etale covering} of~$U$.
\end{theorem}

The ``only if'' being trivial, let us prove the ``if''.
Let $z$ be the generic point of~$Y$; one has $n(z)=$
(sum of the separable degrees of the $K_i/K$) $\le n$,
and by~\ref{I.10.7}, $n(y)\le n(z)\le n$, with equality
implying that $X|U$ is \emph{finite} over~$U$ for a
suitable neighborhood~$U$ of~$y$. One may therefore
suppose that $X$ is finite over~$Y$ and that the
function $n(y')$ on~$Y$ is constant. Finally,
by~\ref{I.10.8}, $X$ is then the disjoint union of
its irreducible components, and to prove that it
is unramified at $y$, one is reduced to the case
where $X$ is irreducible, hence integral. Finally,
one may suppose that $Y=\Spec(\cal{O}_y)$. The theorem
then reduces to the following classical statement:

\begin{corollary}
\label{I.10.12}
Let $A$ be a normal local ring (noetherian, as always)
with field $K$, $L$ a finite extension of~$K$ of
degree $n$ and separable degree $n_s$, $B$ a subring
of~$L$ finite over~$A$, with field of fractions $L$,
$\goth{m}$ the maximal ideal of~$A$, and $n'$ the
separable degree of~$B/\goth{m}B$ over~$A/\goth{m}A=k$
($=$ the sum of the separable degrees of the residue
field extensions of this ring). One has $n'\le n_s$
and a fortiori $n'\le n$. This last inequality is
an equality if and only if $B$ is unramified
($=$ \'etale) over~$A$.
\end{corollary}

It remains only to show that $n'=n$ implies that $B$
is \'etale over~$A$. Let us recall the proof when $k$
is infinite: one need only show that $R=B/\goth{m}B$
is separable over~$k$; if this were not so, it would
follow (by a known lemma) that there exists an
element $a$ of~$R$ whose minimal polynomial over~$k$
has degree $>n'$. This element comes from an element
$x$ of~$B$, whose minimal polynomial over~$K$
(as an element of $L$) has degree $\le n$; on the
other hand, the latter has its coefficients in~$A$
since $A$ is normal, and hence gives, by reduction
mod $\goth{m}$, a monic polynomial $F\in k[t]$ of
degree $\le n=n'$ such that $F(a)=0$, a contradiction.

In
% original p. 26
the general case ($k$ possibly finite), returning
to geometric language, one considers $Y'=\Spec(A[t])$,
which is faithfully flat over~$Y$, and the generic
point $y'$ of the fiber $\Spec(k[t])$ of~$Y'$ over~$y$.
Then $X$ is net over~$Y$ at $y$ if and only if
$X'=X\times_Y Y'=\Spec(B[t])$ is net at $y'$ over~$Y'$,
as one sees at once. Moreover, by the choice of~$y'$,
its residue field is $k(t)$, hence infinite. Since
$y'$ is a normal point of~$Y'$, one is reduced to
the preceding case.
